% ======================================================================
% 第 8 章：引子：紫外截断与临界涨落
% （原书第 265-274 页；PDF 第 288-297 页）
% ======================================================================
\chapter{引子：紫外截断与临界涨落}

本书第~II 部分的主要目的是发展重整化的一般理论。这一理论将解释场论中紫外发散的起源，并指出这些发散何时可以被系统地消除。它还将给出一种方法，把费曼图的发散从一个问题转化为一种工具。我们将应用这一工具来研究场论振幅在大动量或小动量极限下的渐近行为。

当我们在第~6.3 节计算单圈顶点修正时第一次遇到紫外发散时，它看起来像是一种畸变，理应在给我们带来太多困扰之前就消失。在第~7 章中，我们看到了更多紫外发散图的例子，这足以使我们相信此类发散在费曼图计算中普遍存在。因此，任何学习场论的人都有必要对这类发散建立一种观点。大多数人一开始都认为，任何含有发散的理论必然是荒谬的。但这种观点过于狭隘，因为它不仅排除了量子场论，甚至排除了点粒子的经典电动力学。

随着经验的积累，人们也许会采取一种更为宽容的态度，即与发散和平共处：只要发散不出现在物理预言中，就可以接受含有发散的理论。在第~7 章中我们看到，电子在重靶上散射的单圈辐射修正中出现的所有发散，都可以通过系统地消去电子的裸质量和裸电荷而改用它们测得的物理值来消除。在第~10 章中，我们将论证 QED 的所有紫外发散，在微扰论的所有阶上，都可以用这种方式消除。因此，只要人们愿意把电子的质量和电荷视为实验测得的参数，QED 微扰论的预言就总是没有发散的。我们还将在第~10 章中证明，QED 属于一类定义良好的场论，在这类理论中，所有紫外发散都可以在从实验取用固定少数几个物理参数后被消除。这些理论被称为\emph{可重整化的}量子场论，它们是仅有的微扰论能给出良好定义预言的理论。

然而，理想的做法是更进一步，尝试从物理上理解为什么会出现发散，以及为什么在一些理论中发散的影响比在另一些理论中更为严重。这种直接处理发散问题的方法是在 20 世纪 60 年代由 Kenneth Wilson 开创的。解决这一问题所需的关键洞见来自 Wilson 等人发现的量子场论与磁体和流体的统计物理之间的对应关系。Wilson 的重整化方法将是第~12 章的主题。本章简要介绍凝聚态物理中为紫外发散问题提供洞见的相关课题。

\subsection*{形式与物理截断}

紫外发散标志着量子场论中计算的量依赖于某个非常大的动量标度，即紫外截断。等价地，在位置空间中，发散的量依赖于某个非常小的距离标度。

在系统连续描述中引入小距离截断的想法在经典场论中同样出现。通常，截断位于原子距离的标度上，在此标度上连续描述不再适用。然而，截断的大小会在连续理论的某些参数中体现出来。例如，在流体动力学中，粘滞系数和声速等参数的大小，恰好是把典型原子半径与速度结合起来所预期的量级。类似地，在磁体中，磁化率可以通过假设翻转一个电子自旋的能量代价约为十分之一电子伏来估算，这与我们从原子物理中所预期的相符。这些系统中的每一个都在原子的标度上拥有一个自然的紫外截断；通过理解原子标度上的物理，我们就能计算决定更大标度上物理的参数。

然而，在量子场论中，我们对非常短距离标度上的基本物理没有精确的认识。因此，我们只能测量诸如电子的物理电荷和质量之类的参数，而不能从第一性原理出发计算它们。在这些物理参数与其裸值之间的关系中出现紫外发散，正是这些参数受未知短距离物理控制的标志。

无论我们是否知道小距离标度上的基本物理，要为长距离现象写出一个有效理论，我们都需要两类信息。第一，我们必须知道小距离标度中有多少参数对长距离物理是相关的。第二，也是更重要的，我们必须知道底层理论中的哪些自由度会出现在长距离上。

在流体力学中，从原子的观点看，居然存在任何长距离自由度，这多少有点奇迹的意味。尽管如此，描述能量和质量在长距离上输运的方程确实具有光滑、相干的解。长距离自由度是输运这些守恒量的流动，以及长波长的声波。

在量子场论中，长距离物理只涉及那些质量与基本截断标度相比非常小的粒子。这些粒子及其动力学是流体力学中大尺度流动的量子类比。自然地安排此类粒子出现的最简单方法是利用天然具有零质量的粒子。到目前为止，我们在本书中遇到了两类质量精确为零的粒子：光子和手征费米子。（在第~11 章中，我们还将遇到一种天然无质量的粒子，即戈德斯通玻色子。）我们可以论证，QED 之所以能作为远大于其截断标度上的一种理论而存在，是因为光子天然无质量，并且因为左右手电子非常接近手征费米子。

在量子场论中，还有另一种方式可以让零质量或近似零质量的粒子出现：我们可以简单地调节标量场论的参数，使标量粒子的质量远小于截断标度。这种引入小质量粒子的方法看起来有些任意而不自然。尽管如此，它在统计力学中有一个在该学科中真正令人感兴趣的类比，并且能教会我们一些重要的道理。

通常，在凝聚态系统中，热涨落只在原子距离上关联。然而，在特殊情况下，它们可以具有长得多的关联程。这一现象最清晰的例子出现在铁磁体中。在高温下，磁体中的电子自旋是无序的并涨落不定；但在低温下，这些自旋排列到一个固定的方向。\footnote{在真实的铁磁体中，长程的磁偶极—偶极相互作用会使均匀磁化的状态分裂为一系列磁畴。在本书中，我们将忽略这种相互作用，把磁性自旋视为纯粹的取向。正是这个理想化的系统与量子场论直接类似。}让我们考虑随着磁体温度降低，这种排列是如何逐渐形成的。当磁体从高温冷却时，关联自旋的团簇变得越来越大。在某个时刻——磁化温度——整个样品变成一个具有明确宏观取向的单一团簇。就在这个温度之上，磁体中含有具有共同取向的大团簇自旋，而这些团簇又隶属于更大的团簇，使得最大尺度上的取向在样品中仍然是随机化的。这种情况如图~\ref{fig:8.1} 所示。类似的行为也出现在任何其他二级相变的邻域，例如二元合金中的有序—无序转变、流体中的临界点，或氦-4 中的超流转变。

\begin{figure}[t]
  \centering
  \includegraphics[width=0.55\textwidth]{fig8.1.png}
  \caption{铁磁体临界点附近定向自旋的团簇。}
  \label{fig:8.1}
\end{figure}

对这些极长波长涨落的自然描述是借助一个涨落的连续场。在最直观的层面上，我们或许可以用量子涨落取代统计涨落，并尝试把这个系统描述为一种量子场论。在第~9.3 节中，我们将推导一个更为微妙的关系，在统计系统与量子系统之间建立精确的联系。通过这种联系，任何统计系统在二级相变附近的行为都可以转译为某一特定量子场论的行为。这个量子场论包含一个场，其质量与基本原子标度相比非常小，并且在相变点精确地趋于零。

但这种联系似乎加重了量子场论中紫外发散的问题：如果自然界中观察到的丰富多样的相变产生了同样丰富多样的量子场论，那么我们怎么可能在不详细参照其紫外截断标度上的物理起源的情况下定义一个量子场论呢？说一个量子场论的预言与截断无关，就相当于说临界点邻域的统计涨落与该系统是磁体、流体还是合金无关。但是，这种说法真的明显错误吗？通过反转逻辑，我们会发现量子场论对凝聚态系统作出一个非常有力的预言，即对临界点附近统计涨落的\emph{普适性}预言。事实上，这个预言已经得到了实验验证。

本书第~II 部分的一个主要论题是，这两个想法——量子场论中的截断无关性和临界现象理论中的普适性——在本质上是同一个想法，理解其中任何一个都会使我们对另一个获得洞见。

\subsection*{朗道相变理论}

为了对相变现象中什么是普适的有一个初步概念，让我们考察由 Landau 提出的最简单的二级相变连续理论。

首先我们应该回顾一点热力学并澄清我们的术语。在热力学中，\emph{一级相变}是某个热力学变量（流体的密度，或铁磁体的磁化强度）发生不连续变化的点。在相变点，两个截然不同的热力学状态（液体与气体，或平行与反平行于给定轴的磁化）处于平衡。横跨相变不连续变化、并刻画两个竞争相差异的热力学量称为\emph{序参量}。在大多数情况下，可以改变第二个热力学参数，使得两个竞争状态在热力学空间中相互靠近，从而在这个参数的某个值处，这两个状态变得相同，序参量的不连续性消失。这个一级相变线的端点被称为\emph{二级相变}，或者更确切地说，\emph{临界点}。从另一个方向看，临界点是一个单一热力学状态分叉为两个宏观上不同状态的点。正是这种分叉导致了前一节讨论的长程热涨落。

铁磁体展示了一个具体例子。为简单起见，让我们假设所讨论的材料有一个优先磁化轴，因此在低温下，系统的自旋将要么平行、要么反平行于该轴有序排列。沿此轴的总磁化强度 $M$ 就是序参量。在低温下，施加外磁场 $H$ 将有利于两种可能状态中的这一种或那一种。在 $H = 0$ 时，两种状态处于平衡；如果 $H$ 从小的负值变为小的正值，热力学状态和 $M$ 的值将发生不连续变化。因此，对于任何固定的（低温）温度，在 $H = 0$ 处都存在一级相变。现在考虑升高温度的影响：自旋的涨落增大，$|M|$ 的值减小。在某个温度 $T_c$ 处，系统在 $H = 0$ 时不再被磁化。在这个点上，一级相变消失，两个竞争的热力学状态合并。因此，该系统在 $T = T_c$ 处有一个临界点。这些不同的相变在 $H$--$T$ 平面上的位置如图~\ref{fig:8.2} 所示。

\begin{figure}[t]
  \centering
  \includegraphics[width=0.5\textwidth]{fig8.2.png}
  \caption{单轴铁磁体在 $H$--$T$ 平面上的相图。}
  \label{fig:8.2}
\end{figure}

Landau 借助吉布斯自由能 $G$ 描述了这种行为；这是一个依赖于 $M$ 和 $T$ 的热力学势，使得
\begin{equation}
  H = \frac{\partial G}{\partial M}\Big|_{T} .
  \label{eq:8.1}
\end{equation}
他建议我们把注意力集中在临界点的区域：$T\approx T_c$，$M\approx 0$。于是，把 $G(M)$ 展开为 $M$ 的泰勒级数是合理的。对于 $H = 0$，我们可以写出
\begin{equation}
  G(M) = A(T) + B(T)M^2 + C(T)M^4 + \cdots .
  \label{eq:8.2}
\end{equation}
由于系统在 $M\to -M$ 下具有对称性，$G(M)$ 只能包含 $M$ 的偶次幂。由于 $M$ 很小，我们将忽略展开式中的高次项。给定式~\eqref{eq:8.2}，我们可以通过求解
\begin{equation}
  0 = \frac{\partial G}{\partial M} = 2B(T)M + 4C(T)M^3 .
  \label{eq:8.3}
\end{equation}
来找出 $H = 0$ 时 $M$ 的可能取值。如果 $B$ 和 $C$ 都为正，唯一的解是 $M = 0$。然而，如果 $C > 0$ 但 $B$ 在某个温度 $T_c$ 以下为负，则对 $T < T_c$ 存在非平凡解，如图~\ref{fig:8.3} 所示。更具体地说，在 $T\approx T_c$ 附近近似：
\begin{equation}
  B(T) = b(T-T_c),\qquad C(T) = c .
  \label{eq:8.4}
\end{equation}
那么式~\eqref{eq:8.3} 的解为
\begin{equation}
  M \sim \begin{cases}
    0 & \text{for } T > T_c;\\
    \pm \bigl\{ (b/2c)(T_c - T) \bigr\}^{1/2} & \text{for } T < T_c .
  \end{cases}
  \label{eq:8.5}
\end{equation}
这正是我们在临界点所预期的定性行为。

\begin{figure}[t]
  \centering
  \includegraphics[width=0.5\textwidth]{fig8.3.png}
  \caption{朗道理论中吉布斯自由能 $G(M)$ 在临界温度以上和以下的行为。}
  \label{fig:8.3}
\end{figure}

为了求非零外场下 $M$ 的值，我们可以求解式~\eqref{eq:8.1}，其左边由 \eqref{eq:8.2} 给出。一个等价的做法是最小化一个与 \eqref{eq:8.2} 相关的新函数。定义
\begin{equation}
  G(M, H) = A(T) + B(T)M^2 + C(T)M^4 - HM .
  \label{eq:8.6}
\end{equation}
那么，在固定 $H$ 下 $G(M,H)$ 对 $M$ 的最小值给出满足式~\eqref{eq:8.1} 的 $M$ 值。除 $H = 0$ 且 $T < T_c$ 的情形外，该最小值是唯一的，此时我们在 \eqref{eq:8.5} 的第二行看到双重最小值。这与图~\ref{fig:8.2} 所示的相图一致。

为了研究相变邻域的关联，Landau 进一步推广了这一描述，把磁化强度 $M$ 视为局域自旋密度的积分：
\begin{equation}
  M = \int d^3x\ s(\mathbf{x}) .
  \label{eq:8.7}
\end{equation}
于是，吉布斯自由能 \eqref{eq:8.6} 变成 $s(\mathbf{x})$ 的局域函数的积分，
\begin{equation}
  G = \int d^3x\ \Bigl[ \tfrac{1}{2}(\nabla s)^2 + b(T-T_c)s^2 + cs^4 - Hs
    \Bigr],
  \label{eq:8.8}
\end{equation}
它必须对场位形 $s(\mathbf{x})$ 取极小。第一项是引入邻近自旋彼此排列趋向的最简单方式。我们已经对 $s(\mathbf{x})$ 重新标度，使得这一项的系数设为 $1/2$。在写出这个自由能积分时，我们甚至可以认为 $H$ 随位置变化。事实上，这样做是有用的；我们可以在 $\mathbf{x} = 0$ 附近开启 $H(\mathbf{x})$，并观察在另一点会得到什么响应。

自由能表达式 \eqref{eq:8.8} 对 $s(\mathbf{x})$ 的极小值由变分方程
\begin{equation}
  0 = \frac{\delta G[s]}{\delta s(\mathbf{x})} = -\nabla^2 s + 2b(T-T_c)s
    + 4cs^3 - H(\mathbf{x}) .
  \label{eq:8.9}
\end{equation}
的解给出。对于 $T > T_c$，宏观磁化强度为零，因此 $s(\mathbf{x})$ 应该很小，我们可以通过忽略 $s^3$ 项来找出定性行为。于是 $s(\mathbf{x})$ 服从线性方程，
\begin{equation}
  \bigl( -\nabla^2 + 2b(T-T_c) \bigr)\,s(\mathbf{x}) = H(\mathbf{x}) .
  \label{eq:8.10}
\end{equation}
为了研究自旋的关联，我们将令
\begin{equation}
  H(\mathbf{x}) = H_0\,\delta^{(3)}(\mathbf{x}) .
  \label{eq:8.11}
\end{equation}
由此得到的位形 $s(\mathbf{x})$ 正是式~\eqref{eq:8.10} 中微分算符的格林函数，因此我们把它记为 $D(\mathbf{x})$：
\begin{equation}
  \bigl( -\nabla^2 + 2b(T-T_c) \bigr)\,D(\mathbf{x}) = H_0\,\delta^{(3)}(\mathbf{x}).
  \label{eq:8.12}
\end{equation}
这个格林函数告诉我们，当 $\mathbf{x} = 0$ 处的自旋被强制与 $H$ 排列时，在 $\mathbf{x}$ 处会有什么响应。在第~9.2 和~9.3 节中，我们将看到 $D(\mathbf{x})$ 还正比于热系综中零场的自旋—自旋关联函数，
\begin{equation}
  D(\mathbf{x}) \propto \langle s(\mathbf{x})s(0)\rangle
  = \sum_{\text{all } s(\mathbf{x})} s(\mathbf{x})\,s(0)\,
    e^{-\mathcal{H}/k_BT} ,
  \label{eq:8.13}
\end{equation}
其中 $\mathcal{H}$ 是磁系统的哈密顿量。

式~\eqref{eq:8.12} 的解可以通过傅里叶变换求得：
\begin{equation}
  D(\mathbf{x}) = \int \frac{d^3k}{(2\pi)^3}\
    \frac{H_0\,e^{i\mathbf{k}\cdot\mathbf{x}}}{|\mathbf{k}|^2
      + 2b(T-T_c)} .
  \label{eq:8.14}
\end{equation}
这正是我们在讨论汤川势时遇到的积分，即式~(4.126)。用同样的方法计算它，我们得到
\begin{equation}
  D(\mathbf{x}) = \frac{H_0}{4\pi r}\ e^{-r/\xi},
  \label{eq:8.15}
\end{equation}
其中
\begin{equation}
  \xi = \bigl[ 2b(T-T_c) \bigr]^{-1/2}
  \label{eq:8.16}
\end{equation}
是\emph{关联长度}，即关联自旋涨落的程。注意，这个长度在 $T\to T_c$ 时发散。

这一分析的主要结果，即式~\eqref{eq:8.5} 和 \eqref{eq:8.16}，涉及依赖于原子标度物理的未知常数 $b$、$c$。另一方面，这些公式中对 $(T-T_c)$ 的幂律依赖仅仅由朗道方程的结构推出，与微观物理的任何细节无关。事实上，我们在推导这一依赖关系时甚至没有用到 $G$ 描述的是铁磁体这一事实；我们只假设 $G$ 可以按序参量的幂次展开，并且 $G$ 满足反射对称性 $M\to-M$。这些假设同样适用于许多其他类型的系统：二元合金、超流体，甚至（尽管这里的反射对称性不那么明显）液—气相变。朗道理论预言，关联长度的行为为
\begin{equation}
  \xi \sim (T-T_c)^{-\nu},
  \label{eq:8.17}
\end{equation}
在所有这些系统中具有相同的指数 $\nu = 1/2$。这个普适预言是我们在临界现象理论中更一般地预期的一个简单例子。

然而，在真实的物理系统中，$\nu = 1/2$ 的预言并不完全正确。对于三维单轴磁体，观测值为 $\nu = 0.630$；对于流体，$\nu = 0.627$。对于给定\emph{普适类}中的所有系统，这些值都是相同的（在实验误差范围内），其中普适类由空间的维数、序参量的分量数以及系统在序参量反射对称性下的行为所确定。对于流体，相关的序参量是密度，它在反射下是一个单分量标量，因此流体与单轴铁磁体属于同一个普适类。

让我们尝试理解朗道预言与临界指数观测值之间的差异。我们在推导式~\eqref{eq:8.17} 时忽略了热涨落的影响。但恰好位于临界点时，关联长度发散，因此涨落会在任意长的距离上关联，我们不应把它们当作小扰动来处理。对这些涨落的正确处理是本书其余部分的主题。尽管如此，朗道理论仍是一个有用的出发点，因为它正确地预言了临界指数的普适性和关联长度的存在；它只是未能预言精确的数值。

为了改进 Landau 的估计，我们需要系统地考虑涨落。最自然的做法是用连续场 $s(\mathbf{x})$ 写出统计配分函数：
\begin{equation}
  Z = \int \mathcal{D}s(\mathbf{x})\ \exp\Bigl[ -\frac{1}{k_BT}\int d^3x\
    \bigl( \tfrac{1}{2}(\nabla s)^2 + b(T-T_c)s^2 + cs^4 - Hs \bigr) \Bigr].
  \label{eq:8.18}
\end{equation}
这是对涨落场 $s(\mathbf{x})$ 的泛函积分。这一对象的出现标志着统计力学与量子场论之间的深刻联系：泛函积分 \eqref{eq:8.18} 正是我们在第~9 章中用来定义量子场论的那类对象。通过把温度等同于虚时间的倒数，这种对应关系被精确化，我们将在第~9.3 节讨论这一点。

涨落会改变幂律。例如，式~\eqref{eq:8.16} 在 $T \to T_c$ 时变为
\begin{equation}
  \xi \sim (T-T_c)^{-\nu},
  \label{eq:8.19}
\end{equation}
其中指数 $\nu$ 对给定普适类中的所有系统取一个固定值。对于三维单轴磁体和流体，$\nu = 0.630$。这些非平凡标度关系中的幂次被称为\emph{临界指数}。

从式~\eqref{eq:8.18} 到式~\eqref{eq:8.19} 的修改并不会动摇这一观念，即处于二级相变邻域的凝聚态系统具有定义良好、与截断无关的连续行为。然而，我们希望理解为什么式~\eqref{eq:8.17} 应当被预期为正确的表述。这一问题的答案将来自对相应量子场论紫外发散的透彻分析。在第~12 章中，当我们最终完成对紫外发散的阐述时，我们将发现我们手中不仅有证明式~\eqref{eq:8.17} 的工具，还有用费曼图计算临界指数数值的工具。通过这种方式，我们将揭示量子场论在原子物理领域的一个优美应用。这一应用的成功将在第~III 部分引导我们获得更强大的工具，这些工具将在基本粒子的相对论性领域为我们所需要。
